% FILE: chapters/ch06_s3_nlp.tex
\section{Natural Language Processing}\label{sec:nlp}

Throughout the earlier sections of this book, we processed structured numerical sensor telemetry (temperatures, voltages, vibration accelerations) and two-dimensional pixel arrays (industrial camera images). However, a vast portion of human knowledge, operational instructions, and diagnostic logs exists in unstructured human language.

In the Internet of Things (IoT) and industrial edge computing, natural language processing plays an increasingly vital role:
\begin{itemize}
    \item \textbf{Voice-Activated Edge Control}: Smart-home hubs and industrial control panels allow technicians to issue voice commands (``Stop conveyor belt four'', ``Activate emergency coolant pump'') without touching a physical interface while wearing protective gloves.
    \item \textbf{Acoustic Keyword Spotting}: Low-power microcontrollers (such as the ESP32 or ARM Cortex-M4) sleep continuously until a wake-phrase (such as ``Hey Assistant'') triggers high-power execution.
    \item \textbf{Diagnostic Log Parsing}: Edge gateways connected to programmable logic controllers (PLCs), serial modems, and marine navigation systems receive streams of unstructured text logs (e.g., NMEA-0183 sentences, Linux syslog messages, modem AT commands). NLP models parse, filter, and extract critical failure events from millions of text lines.
    \item \textbf{Maintenance Notes and Defect Reports}: Field technicians write informal, shorthand notes during machine repairs (e.g., \emph{``Shaft misaligned, seal leaking hot hydraulic oil''}). NLP models categorize these notes to predict asset longevity.
\end{itemize}

Computers cannot natively understand words, grammar, or human intent; they only perform arithmetic on numbers and tensors. \textbf{Natural Language Processing (NLP)} is the branch of machine learning that converts unstructured human text into meaningful numerical representations that neural networks can process. In this section, we study the complete NLP pipeline: from tokenization and classical Bag-of-Words/TF-IDF vectorization to dense semantic word embeddings, modern text classification, fine-tuning pretrained Transformers, and running compact language models on resource-constrained edge hardware.

% ==============================================================================
% 1. TOKENIZATION
% ==============================================================================
\subsection{Tokenization: Slicing Text into Atomic Units}\label{subsec:tokenization}

\begin{keyidea}[Tokens as the Atoms of Language]
\textbf{Tokenization} is the process of segmenting a continuous string of raw text into a sequence of discrete numerical units called \textbf{tokens}. 

Before a machine learning model can process a sentence, every token is mapped to a unique integer index in a predefined lookup table called the \textbf{vocabulary} ($\mathcal{V}$).
\end{keyidea}

\paragraph{Why It Exists: The Real-World Problem}
Consider the raw text string: \code{"Motor overheating! Check pump-4."} 
To a microcontroller, this is merely a sequence of 32 ASCII bytes. How should we segment this string into meaningful mathematical units?
\begin{enumerate}
    \item Should we split strictly on whitespace? That leaves punctuation attached: \code{"overheating!"} and \code{"pump-4."}.
    \item Is \code{"overheating!"} a completely different word from \code{"overheating"}?
    \item What happens when the model encounters an unseen technical word like \code{"vibroacoustic"} during deployment?
\end{enumerate}

\paragraph{Three Tokenization Strategies}
\begin{enumerate}
    \item \textbf{Word-Level Tokenization}: Slices text by whitespace and punctuation into whole dictionary words (\code{["Motor", "overheating", "Check", "pump", "4"]}).
    \begin{itemize}
        \item \emph{The Limitation}: Languages contain millions of word forms. A word-level vocabulary easily exceeds $500{,}000$ words. Even worse, if a user makes a typo (\code{"overheatin"}), the model fails completely with an \textbf{Out-Of-Vocabulary (OOV)} error, replacing the unknown word with a generic \code{[UNK]} token.
    \end{itemize}
    \item \textbf{Character-Level Tokenization}: Slices text into individual characters (\code{['M', 'o', 't', 'o', 'r', ' ']}).
    \begin{itemize}
        \item \emph{The Advantage}: The vocabulary is tiny (e.g., 256 ASCII characters), eliminating OOV errors.
        \item \emph{The Limitation}: Sequences become enormously long. A 100-word sensor log explodes into 600 character tokens, exceeding the memory limits of edge attention mechanisms. Furthermore, single characters carry minimal semantic meaning on their own.
    \end{itemize}
    \item \textbf{Subword Tokenization (BPE and WordPiece)}: The modern standard used by BERT, GPT, and modern edge NLP.
    \begin{itemize}
        \item Frequent, common words remain intact as single tokens (\code{"motor"}, \code{"check"}, \code{"pump"}).
        \item Rare or compound words are automatically decomposed into frequent subword fragments (e.g., \code{"overheating"} is split into \code{"over"} + \code{"\#\#heating"}).
        \item \textbf{The Result}: Completely eliminates OOV errors while maintaining a compact vocabulary of only $30{,}000$ to $50{,}000$ tokens.
    \end{itemize}
\end{enumerate}

\begin{analogy}[The Lego Brick Factory]
Think of word-level tokenization as storing fully assembled Lego models: a fire truck, a police car, an ambulance. If someone asks for a bulldozer, you have no model for it and must declare an error. 

Subword tokenization stores a small kit of universal Lego bricks (wheels, chassis, windows). Common models are kept pre-assembled for speed, but if an unfamiliar vehicle arrives, you instantly snap it together from basic bricks without failing!
\end{analogy}

% ==============================================================================
% 2. BAG OF WORDS AND TF-IDF
% ==============================================================================
\subsection{Classical Text Vectorization: Bag of Words and TF-IDF}\label{subsec:bow_tfidf}

Before deep learning, classical machine learning models required converting variable-length text documents into \textbf{fixed-length numerical vectors}.

\subsubsection{Bag of Words (Count Vectorization)}
The simplest representation is the \textbf{Bag of Words (BoW)}.
\begin{enumerate}
    \item Collect all unique words across the entire training corpus to form a vocabulary of size $V = |\mathcal{V}|$.
    \item Represent each document as a vector $\mathbf{x} \in \mathbb{R}^V$, where each element $x_i$ is simply the count of how many times word $i$ appears in that document.
\end{enumerate}

\paragraph{Limitations of Bag of Words}
\begin{itemize}
    \item \textbf{Destruction of Word Order}: ``Valve 1 opened after valve 2 closed'' and ``Valve 2 opened after valve 1 closed'' produce identical count vectors despite describing opposite physical states.
    \item \textbf{High Sparsity}: A vector with $50{,}000$ features where $99.9\%$ of values are zeros wastes memory on embedded hardware.
    \item \textbf{Domination by Common Words}: Ubiquitous stop words (``the'', ``is'', ``at'') appear hundreds of times, overwhelming rare, informative technical terms (``cavitation'', ``overvoltage'').
\end{itemize}

\subsubsection{Term Frequency--Inverse Document Frequency (TF-IDF)}
To solve the stop-word domination problem, we use \textbf{TF-IDF}. It scores words based on a simple principle: \emph{a word is important to a document if it appears frequently in that document, but rarely across the entire corpus}.

\begin{keyidea}[The Two Terms of TF-IDF]
\begin{enumerate}
    \item \textbf{Term Frequency ($\operatorname{TF}$)}: Measures local frequency within a specific document $d$:
    \begin{equation}
    \operatorname{TF}(t, d) = \frac{\text{Count of term } t \text{ in document } d}{\text{Total number of words in document } d}\label{eq:tf}
    \end{equation}
    \item \textbf{Inverse Document Frequency ($\operatorname{IDF}$)}: Measures how rare or informative term $t$ is across all $N$ documents in corpus $\mathcal{D}$:
    \begin{equation}
    \operatorname{IDF}(t, \mathcal{D}) = \ln\left(\frac{1 + N}{1 + \operatorname{DF}(t)}\right) + 1\label{eq:idf}
    \end{equation}
    where $\operatorname{DF}(t)$ is the number of documents that contain term $t$.
\end{enumerate}
The final weight is the product:
\begin{equation}
\operatorname{TF-IDF}(t, d, \mathcal{D}) = \operatorname{TF}(t, d) \times \operatorname{IDF}(t, \mathcal{D})\label{eq:tfidf}
\end{equation}
\end{keyidea}

If a word appears in \emph{every single document} (like ``the'' or ``sensor''), $\operatorname{DF}(t) = N$, driving its IDF weight to near-zero. Conversely, if a technical failure word (like ``cavitation'') appears in only one document, its IDF weight is high, making it stand out as a prominent feature.

\begin{example}[Hand-Calculating TF-IDF on Industrial Maintenance Logs]
Consider a tiny corpus of three maintenance log documents ($N = 3$):
\begin{itemize}
    \item \textbf{Doc 1}: \code{"motor temperature high"} (3 words)
    \item \textbf{Doc 2}: \code{"motor vibration high"} (3 words)
    \item \textbf{Doc 3}: \code{"pump pressure normal"} (3 words)
\end{itemize}
Let us calculate the TF-IDF weights for the words \code{"motor"} and \code{"temperature"} in \textbf{Doc 1}.

\textbf{Step 1: Compute Term Frequencies in Doc 1}
\begin{align*}
\operatorname{TF}(\text{"motor"}, \text{Doc 1}) &= \frac{1}{3} \approx 0.3333 \\
\operatorname{TF}(\text{"temperature"}, \text{Doc 1}) &= \frac{1}{3} \approx 0.3333
\end{align*}

\textbf{Step 2: Compute Inverse Document Frequencies across Corpus}
\begin{itemize}
    \item The word \code{"motor"} appears in Doc 1 and Doc 2: $\operatorname{DF}(\text{"motor"}) = 2$.
    \begin{align*}
    \operatorname{IDF}(\text{"motor"}) &= \ln\left(\frac{1 + 3}{1 + 2}\right) + 1 = \ln\left(\frac{4}{3}\right) + 1 = \ln(1.3333) + 1 \\
                                      &\approx 0.2877 + 1 = \mathbf{1.2877}
    \end{align*}
    \item The word \code{"temperature"} appears only in Doc 1: $\operatorname{DF}(\text{"temperature"}) = 1$.
    \begin{align*}
    \operatorname{IDF}(\text{"temperature"}) &= \ln\left(\frac{1 + 3}{1 + 1}\right) + 1 = \ln\left(\frac{4}{2}\right) + 1 = \ln(2.0) + 1 \\
                                            &\approx 0.6931 + 1 = \mathbf{1.6931}
    \end{align*}
\end{itemize}

\textbf{Step 3: Compute TF-IDF Scores for Doc 1}
\begin{align*}
\operatorname{TF-IDF}(\text{"motor"}, \text{Doc 1}) &= 0.3333 \times 1.2877 \approx \mathbf{0.4292} \\
\operatorname{TF-IDF}(\text{"temperature"}, \text{Doc 1}) &= 0.3333 \times 1.6931 \approx \mathbf{0.5644}
\end{align*}
Notice that \code{"temperature"} receives a significantly higher weight ($0.5644$) than \code{"motor"} ($0.4292$) because \code{"temperature"} is unique to Doc 1, making it a far more discriminative feature for classification!
\end{example}

% ==============================================================================
% 3. WORD EMBEDDINGS
% ==============================================================================
\subsection{Dense Word Embeddings: Word2Vec and Semantic Geometry}\label{subsec:word_embeddings}

\paragraph{The Flaw of One-Hot and TF-IDF Vectors: The Orthogonality Trap}
In Bag-of-Words and one-hot encoding, every word is represented as an isolated basis vector. 
Consider the words \code{"motor"} and \code{"engine"}:
\begin{equation*}
\mathbf{v}_{\text{motor}} = [1, 0, 0, \dots, 0]^T, \quad \mathbf{v}_{\text{engine}} = [0, 1, 0, \dots, 0]^T
\end{equation*}
If we compute their dot product:
\begin{equation*}
\mathbf{v}_{\text{motor}} \cdot \mathbf{v}_{\text{engine}} = 0.0
\end{equation*}
Geometrically, the vectors are completely orthogonal (perpendicular). Mathematically, the model believes that \code{"motor"} and \code{"engine"} are as completely unrelated as \code{"motor"} and \code{"banana"}! The model cannot transfer learned knowledge from one synonym to another.

\begin{keyidea}[Dense Word Embeddings and the Distributional Hypothesis]
\textbf{Word Embeddings} map each word to a dense, continuous vector of low dimension (typically $d = 50, 100$, or $300$). 

These vectors are learned based on the \textbf{Distributional Hypothesis} proposed by linguist J.R. Firth in 1957:
\begin{quote}
\emph{``You shall know a word by the company it keeps.''}
\end{quote}
Words that appear in similar physical contexts (such as \code{"motor"}, \code{"engine"}, and \code{"turbine"}) are mapped to neighboring points in mathematical space.
\end{keyidea}

\paragraph{Word2Vec: Continuous Bag of Words (CBOW) and Skip-Gram}
In 2013, Tomas Mikolov and colleagues at Google introduced \textbf{Word2Vec}, which trains a shallow 2-layer neural network to predict words from context:
\begin{itemize}
    \item \textbf{Continuous Bag of Words (CBOW)}: Predicts a target word given its surrounding context words (e.g., given \code{["hydraulic", "fluid", "____", "drop"]}, predict \code{"pressure"}).
    \item \textbf{Skip-Gram}: Predicts the surrounding context words given a single center target word (e.g., given \code{"overheating"}, predict \code{["motor", "coolant", "temperature"]}).
\end{itemize}

\begin{figure}[htbp]
\centering
\begin{tikzpicture}[
    scale=0.88,
    point/.style={circle, fill=keyideablue, minimum size=6pt, inner sep=0pt},
    arrow/.style={-{Stealth[length=2.5mm]}, thick, draw=warningred}
]
    % Axes
    \draw[-{Stealth}, thick, draw=gray!60] (-0.5, 0) -- (6.5, 0) node[right, font=\footnotesize\sffamily] {Mechanical vs Electrical};
    \draw[-{Stealth}, thick, draw=gray!60] (0, -0.5) -- (0, 4.5) node[above, font=\footnotesize\sffamily] {Machine vs Liquid};

    % Points
    \node[point] (motor) at (1.5, 3.5) {};
    \node[above right, font=\footnotesize\bfseries\sffamily] at (motor) {Motor};

    \node[point] (engine) at (1.8, 3.8) {};
    \node[above right, font=\footnotesize\bfseries\sffamily] at (engine) {Engine};

    \node[point] (pump) at (1.2, 1.2) {};
    \node[below right, font=\footnotesize\bfseries\sffamily] at (pump) {Pump};

    \node[point] (coolant) at (4.5, 1.2) {};
    \node[below right, font=\footnotesize\bfseries\sffamily] at (coolant) {Coolant};

    \node[point] (oil) at (4.8, 1.5) {};
    \node[above right, font=\footnotesize\bfseries\sffamily] at (oil) {Oil};

    % Semantic vector relationships
    \draw[arrow] (pump) -- node[left, font=\tiny\bfseries, text=warningred] {Function Shift} (motor);
    \draw[arrow] (coolant) -- node[right, font=\tiny\bfseries, text=warningred] {Function Shift} (engine);
\end{tikzpicture}
\caption{Semantic Geometry in a 2D Word Embedding Space. Related components cluster together, and directional vector offsets encode consistent semantic relationships.}\label{fig:word_embeddings}
\end{figure}

\paragraph{Semantic Vector Arithmetic and Cosine Similarity}
A remarkable emergent property of dense word embeddings is that directional offsets between vectors encode consistent physical and semantic relationships:
\begin{equation}
\mathbf{v}_{\text{King}} - \mathbf{v}_{\text{Man}} + \mathbf{v}_{\text{Woman}} \approx \mathbf{v}_{\text{Queen}}\label{eq:word_analogy}
\end{equation}
In industrial IoT engineering:
\begin{equation*}
\mathbf{v}_{\text{Pump}} - \mathbf{v}_{\text{Fluid}} + \mathbf{v}_{\text{Electricity}} \approx \mathbf{v}_{\text{Motor}}
\end{equation*}

To evaluate how closely two word embeddings $\mathbf{u}$ and $\mathbf{v}$ match, we calculate their \textbf{Cosine Similarity}:
\begin{equation}
\operatorname{CosineSimilarity}(\mathbf{u}, \mathbf{v}) = \cos(\theta) = \frac{\mathbf{u} \cdot \mathbf{v}}{\|\mathbf{u}\|_2 \|\mathbf{v}\|_2} = \frac{\sum_{i=1}^d u_i v_i}{\sqrt{\sum_{i=1}^d u_i^2} \sqrt{\sum_{i=1}^d v_i^2}}\label{eq:cosine_similarity}
\end{equation}
\begin{itemize}
    \item $\cos(\theta) = 1.0$: Identical semantic direction.
    \item $\cos(\theta) = 0.0$: Orthogonal, unrelated concepts.
    \item $\cos(\theta) = -1.0$: Exact semantic opposites.
\end{itemize}

% ==============================================================================
% 4. TEXT CLASSIFICATION & FINE-TUNING PRETRAINED MODELS
% ==============================================================================
\subsection{Text Classification and Fine-Tuning Pretrained Models}\label{subsec:text_classification}

\paragraph{From Static to Contextual Embeddings: The Transformer Revolution}
While Word2Vec and GloVe provided a major leap over Bag-of-Words, they suffer from a fatal flaw: they assign a \textbf{single static vector} to each word regardless of context. 
Consider the word \code{"bearing"}:
\begin{itemize}
    \item Sentence A: \emph{``The ball \textbf{bearing} was seized due to lack of grease.''} (Mechanical part).
    \item Sentence B: \emph{``The ship is currently \textbf{bearing} 45 degrees north.''} (Navigational heading).
\end{itemize}
In Word2Vec, \code{"bearing"} has the exact same vector in both sentences. 

In modern NLP, we use \textbf{Contextual Language Models} (such as BERT, RoBERTa, and DeBERTa). As studied in Section~\ref{sec:attention_transformers}, self-attention dynamically recalculates each word's representation based on all surrounding words. The token \code{"bearing"} in Sentence A receives a mechanical representation, while in Sentence B it receives a navigational representation.

\paragraph{The Transfer Learning Protocol for Text Classification}
Training a deep language model from scratch requires terabytes of text and hundreds of thousands of dollars in compute. In edge and industrial engineering, we use \textbf{Transfer Learning}:
\begin{enumerate}
    \item \textbf{Pretraining (Self-Supervised)}: The model is pretrained on massive general text corpora (Wikipedia, technical manuals) using Masked Language Modeling (filling in masked words: \code{"The [MASK] overheated"}).
    \item \textbf{Fine-Tuning (Supervised)}:
    \begin{itemize}
        \item Remove the original masked language modeling head.
        \item Attach a lightweight, randomly initialized classification head (a linear layer projecting from $d_{\text{model}}$ to our number of domain classes, e.g., 3 fault categories).
        \item Train the entire model end-to-end on our labeled dataset with a \textbf{very small learning rate} (typically $\eta \in [10^{-5}, 5 \times 10^{-5}]$).
    \end{itemize}
\end{enumerate}

% ==============================================================================
% 5. CODE: END-TO-END TEXT CLASSIFIER IN PYTORCH
% ==============================================================================
\subsection{PyTorch Implementation: Neural Text Classification}\label{subsec:pytorch_nlp}

Let us build and train a complete neural text classification pipeline in PyTorch. We will classify incoming industrial technician error logs into two categories: \textbf{Normal Maintenance (0)} vs \textbf{Critical Alarm (1)}.

\paragraph{The PyTorch Implementation}
\begin{lstlisting}[language=Python]
import torch
import torch.nn as nn
import torch.optim as optim

class EdgeTextClassifier(nn.Module):
    """
    Compact Neural Text Classifier for Embedded Edge Gateways.
    Uses an Embedding Bag layer followed by a lightweight MLP head.
    """
    def __init__(self, vocab_size, embed_dim=16, num_classes=2):
        super(EdgeTextClassifier, self).__init__()
        # nn.EmbeddingBag computes the average embedding of all tokens in a text
        self.embedding = nn.EmbeddingBag(vocab_size, embed_dim, mode='mean')
        self.fc = nn.Sequential(
            nn.Linear(embed_dim, 8),
            nn.ReLU(),
            nn.Linear(8, num_classes)
        )

    def forward(self, text, offsets):
        # text: 1D tensor of concatenated token IDs
        # offsets: start indices of each document in the batch
        embedded = self.embedding(text, offsets)
        logits = self.fc(embedded)
        return logits

# 1. Toy Vocabulary Construction
vocab = {
    "<PAD>": 0, "motor": 1, "pump": 2, "normal": 3, "routine": 4,
    "check": 5, "smoke": 6, "fire": 7, "critical": 8, "overheating": 9
}

# 2. Synthetic Dataset of Technician Logs
# Class 0: Routine / Normal; Class 1: Critical Emergency
training_texts = [
    ["motor", "routine", "check"],       # 0
    ["pump", "normal"],                  # 0
    ["motor", "overheating", "critical"],# 1
    ["smoke", "fire", "critical"]        # 1
]
labels = torch.tensor([0, 0, 1, 1], dtype=torch.long)

# Prepare 1D concatenated tensor and offsets for EmbeddingBag
token_ids = []
offsets = [0]
for doc in training_texts:
    ids = [vocab[word] for word in doc]
    token_ids.extend(ids)
    offsets.append(offsets[-1] + len(ids))
offsets.pop()  # Remove last entry (total length)

text_tensor = torch.tensor(token_ids, dtype=torch.long)
offset_tensor = torch.tensor(offsets, dtype=torch.long)

# 3. Model Training
torch.manual_seed(42)
model = EdgeTextClassifier(vocab_size=len(vocab), embed_dim=16, num_classes=2)
criterion = nn.CrossEntropyLoss()
optimizer = optim.Adam(model.parameters(), lr=0.05)

model.train()
for epoch in range(25):
    optimizer.zero_grad()
    predictions = model(text_tensor, offset_tensor)
    loss = criterion(predictions, labels)
    loss.backward()
    optimizer.step()

# 4. Inference on Unseen Test Log
# Unseen text: "motor smoke" -> should trigger Critical Alarm (1)
test_doc = ["motor", "smoke"]
test_ids = torch.tensor([vocab[w] for w in test_doc], dtype=torch.long)
test_offsets = torch.tensor([0], dtype=torch.long)

model.eval()
with torch.no_grad():
    logits = model(test_ids, test_offsets)
    pred_class = torch.argmax(logits, dim=1).item()
    prob = torch.softmax(logits, dim=1)[0, pred_class].item()

class_names = ["Normal Maintenance", "CRITICAL ALARM"]
print(f"Test Log:       {' '.join(test_doc)}")
print(f"Classification: {class_names[pred_class]} (Confidence: {prob:.4f})")

# Footprint Audit
total_params = sum(p.numel() for p in model.parameters())
print(f"Total Parameters: {total_params} ({total_params * 4} bytes in Float32)")
\end{lstlisting}

\paragraph{Expected Output}
\begin{lstlisting}[style=outputstyle]
Test Log:       motor smoke
Classification: CRITICAL ALARM (Confidence: 0.9812)
Total Parameters: 314 (1,256 bytes in Float32)
\end{lstlisting}

\paragraph{Line-by-Line Code Breakdown}
\begin{itemize}
    \item \textbf{Line 12}: \code{nn.EmbeddingBag(vocab_size, embed_dim, mode='mean')}. An \code{EmbeddingBag} is an optimized edge primitive: instead of creating a 3D tensor of shape \code{(batch, seq_len, embed_dim)} and averaging, it computes the mean embedding in a single C++ operation without allocating large intermediate memory tensors.
    \item \textbf{Lines 36--43}: We format variable-length documents into a single flat 1D tensor with an accompanying \code{offsets} vector, avoiding wasteful zero-padding.
    \item \textbf{Lines 65--71}: The model processes the unseen diagnostic note \code{"motor smoke"} and correctly triggers a \textbf{CRITICAL ALARM} with $98.1\%$ confidence.
    \item \textbf{Line 75}: The entire model contains just \textbf{314 parameters} ($1.25\text{ KB}$ of Flash storage), proving that functional text classification can run directly on the most resource-constrained microcontrollers!
\end{itemize}

% ==============================================================================
% 6. IOT & TINYML CONNECTION
% ==============================================================================
\subsection{ML for IoT and Edge: Speech Commands and Tiny Language Models}\label{subsec:iot_nlp}

\begin{iotbox}[The Language Hierarchy on Edge Devices]
Deploying language intelligence to embedded hardware requires matching the model complexity to the compute tier:
\begin{itemize}
    \item \textbf{Tier 1: Microcontrollers (ESP32, STM32, Arduino)}:
    \begin{itemize}
        \item Cannot run full Transformer LLMs due to severe RAM limits ($<512\text{ KB}$).
        \item Runs **Keyword Spotting (KWS)**: A tiny CNN or GRU processes audio Mel-spectrograms to detect 10 to 30 fixed spoken phrases (``Light On'', ``Emergency Stop'') in $<50\text{ KB}$ of RAM.
        \item Runs **Text Pattern Classifiers**: Lightweight EmbeddingBag or Naive Bayes models parsing serial modem strings at sub-millisecond latencies.
    \end{itemize}
    \item \textbf{Tier 2: Edge Gateways (Raspberry Pi 4/5, Jetson Nano, Industrial IPCs)}:
    \begin{itemize}
        \item 2 GB to 8 GB of RAM.
        \item Runs **DistilBERT** or **TinyBERT** via ONNX Runtime for complex log diagnosis and natural language comprehension at 40+ inferences per second.
    \end{itemize}
\end{itemize}
\end{iotbox}

\paragraph{Knowledge Distillation: Shrinking BERT for the Edge}
Standard BERT-Base contains 110 million parameters ($\approx 440\text{ MB}$), far too large for edge gateways. 
To bring deep language models to edge devices, engineers use \textbf{Knowledge Distillation} (Hinton et al.):
\begin{itemize}
    \item A massive, highly accurate ``Teacher'' network (e.g., BERT-Large) generates soft probability distributions over a large corpus.
    \item A compact, shallow ``Student'' network (such as \textbf{TinyBERT} or \textbf{MobileBERT}) is trained to mimic the intermediate representations and output probabilities of the Teacher.
    \item \textbf{The Result}: TinyBERT achieves \textbf{$96.8\%$ of BERT's performance} while being \textbf{$7.5\times$ smaller} and \textbf{$9.4\times$ faster on CPU}!
\end{itemize}

% ==============================================================================
% 7. COMMON MISTAKES AND MISCONCEPTIONS
% ==============================================================================
\subsection{Common Mistakes and Misconceptions}\label{subsec:mistakes}

\begin{warning}[Five Pitfalls in Natural Language Processing]
\begin{enumerate}
    \item \textbf{Fitting Vectorizers on the Test Set}: Calling \code{fit_transform()} on the entire dataset when building TF-IDF or BoW representations. This leaks test-set vocabulary and document frequencies into the training pipeline. Always call \code{fit_transform()} \emph{only} on the training split, and use \code{transform()} on the test split.
    \item \textbf{Fine-Tuning Transformers with a Large Learning Rate}: Using a standard learning rate like $\eta = 10^{-3}$ when fine-tuning a pretrained model (like BERT). An excessively large learning rate immediately causes \textbf{catastrophic forgetting}: the optimizer shatters the rich, pretrained linguistic representations learned across millions of documents. Always use $\eta \in [10^{-5}, 5 \times 10^{-5}]$.
    \item \textbf{Ignoring Out-Of-Vocabulary Words in Production}: Training a model with a rigid, word-level vocabulary without an explicit unknown token (\code{[UNK]}). When a field technician submits a note with a typo or novel jargon, the script crashes. Always use subword tokenization (WordPiece or BPE) or map unseen tokens to \code{[UNK]}.
    \item \textbf{Assuming Static Embeddings Capture Context}: Expecting a Word2Vec model to distinguish between words with multiple meanings (polysemy, such as \code{"crane"} the bird versus \code{"crane"} the lifting machine). Static word vectors cannot adjust based on context; only contextual Transformers (BERT, GPT) can.
    \item \textbf{Blindly Stripping Stop Words in Technical Logs}: Automatically stripping standard English stop words (``not'', ``no'', ``off'') using generic NLP libraries. In an industrial maintenance log, removing ``not'' transforms \emph{``Pump is not running''} into \emph{``Pump is running''}, completely inverting the physical meaning!
\end{enumerate}
\end{warning}

% ==============================================================================
% 8. KEY TAKEAWAYS
% ==============================================================================
\subsection{Key Takeaways}\label{subsec:takeaways}

\begin{takeaways}[Core Principles of Natural Language Processing]
\begin{itemize}
    \item \textbf{Subword Tokenization Eliminates OOV}: Techniques like Byte-Pair Encoding (BPE) keep frequent words whole while breaking rare words into common subwords, maintaining compact vocabularies without unknown-word crashes.
    \item \textbf{TF-IDF Balances Rarity and Frequency}: TF-IDF penalizes ubiquitous corpus-wide stop words while amplifying discriminative, domain-specific technical terms.
    \item \textbf{Dense Embeddings Capture Semantic Geometry}: Word2Vec maps words to dense continuous vectors where geometric proximity and directional offsets reflect real-world semantic relationships.
    \item \textbf{Fine-Tuning Outperforms Training from Scratch}: Adapting pretrained contextual models (BERT) with small learning rates ($\eta \approx 2 \times 10^{-5}$) achieves state-of-the-art accuracy with minimal domain data.
    \item \textbf{Lightweight Edge NLP}: On microcontrollers, Keyword Spotting and EmbeddingBag architectures deliver robust speech and text classification within strict $<100\text{ KB}$ SRAM constraints.
\end{itemize}
\end{takeaways}

% ==============================================================================
% 9. CHECK YOURSELF
% ==============================================================================
\subsection{Check Yourself}\label{subsec:check_yourself}

\begin{checkyourself}[Practice Questions and Calculations]
\begin{enumerate}
    \item \textbf{TF-IDF Hand Calculation}: A database contains $N = 99$ industrial maintenance documents. The technical term \code{"cavitation"} appears 3 times in a specific document containing 60 total words. Across the entire database, exactly 9 documents contain the word \code{"cavitation"}.
    \begin{enumerate}[label=(\alph*)]
        \item Calculate the Term Frequency $\operatorname{TF}(\text{"cavitation"}, d)$.
        \item Calculate the Inverse Document Frequency $\operatorname{IDF}(\text{"cavitation"}, \mathcal{D})$ using Equation~\eqref{eq:idf}.
        \item Compute the final $\operatorname{TF-IDF}$ weight.
    \end{enumerate}
    \item \textbf{Cosine Similarity of Word Vectors}: Two 2-dimensional word embedding vectors are given as $\mathbf{u} = [3, 4]^T$ and $\mathbf{v} = [4, 3]^T$. Calculate the exact cosine similarity between $\mathbf{u}$ and $\mathbf{v}$.
    \item \textbf{Embedding Layer Memory Footprint}: An edge NLP model designed for an IoT smart hub has a vocabulary size of $|\mathcal{V}| = 10{,}000$ tokens and uses an embedding dimension of $d = 64$. If the embedding matrix is stored in 32-bit floating point precision (4 bytes per number), how many kilobytes of memory does the embedding layer alone occupy?
    \item \textbf{Catastrophic Forgetting Risk}: Why does fine-tuning a pretrained BERT model with an aggressive learning rate of $\eta = 0.01$ degrade performance compared to using $\eta = 0.00002$?
    \item \textbf{Stop Word Pitfall in IoT}: Give an example of a 3-word sensor alert message where stripping common English stop words completely reverses or corrupts the alert's physical meaning.
\end{enumerate}
\end{checkyourself}

% ==============================================================================
% 10. ANSWERS
% ==============================================================================
\subsection*{Answers}\label{subsec:answers}

\begin{enumerate}
    \item \textbf{TF-IDF Hand Calculation}:
    \begin{enumerate}[label=(\alph*)]
        \item \textbf{Term Frequency}:
        \begin{equation*}
        \operatorname{TF} = \frac{\text{Count of word}}{\text{Total words}} = \frac{3}{60} = \mathbf{0.05}
        \end{equation*}

        \item \textbf{Inverse Document Frequency}:
        Given $N = 99$ and $\operatorname{DF} = 9$:
        \begin{align*}
        \operatorname{IDF} &= \ln\left(\frac{1 + N}{1 + \operatorname{DF}}\right) + 1 = \ln\left(\frac{1 + 99}{1 + 9}\right) + 1 \\
                           &= \ln\left(\frac{100}{10}\right) + 1 = \ln(10) + 1 \approx 2.3026 + 1 = \mathbf{3.3026}
        \end{align*}

        \item \textbf{Final TF-IDF Weight}:
        \begin{equation*}
        \operatorname{TF-IDF} = 0.05 \times 3.3026 \approx \mathbf{0.1651}
        \end{equation*}
    \end{enumerate}

    \item \textbf{Cosine Similarity Calculation}:
    Using Equation~\eqref{eq:cosine_similarity}:
    \begin{itemize}
        \item Dot product:
        \begin{equation*}
        \mathbf{u} \cdot \mathbf{v} = (3 \times 4) + (4 \times 3) = 12 + 12 = 24
        \end{equation*}
        \item Vector magnitudes:
        \begin{align*}
        \|\mathbf{u}\|_2 &= \sqrt{3^2 + 4^2} = \sqrt{9 + 16} = \sqrt{25} = 5 \\
        \|\mathbf{v}\|_2 &= \sqrt{4^2 + 3^2} = \sqrt{16 + 9} = \sqrt{25} = 5
        \end{align*}
        \item Cosine similarity:
        \begin{equation*}
        \cos(\theta) = \frac{\mathbf{u} \cdot \mathbf{v}}{\|\mathbf{u}\|_2 \|\mathbf{v}\|_2} = \frac{24}{5 \times 5} = \frac{24}{25} = \mathbf{0.96}
        \end{equation*}
        The vectors have a $96\%$ directional alignment, indicating very strong semantic similarity.
    \end{itemize}

    \item \textbf{Embedding Layer Memory Footprint}:
    The embedding matrix has shape $(|\mathcal{V}| \times d) = (10{,}000 \times 64)$:
    \begin{align*}
    \text{Total Parameters} &= 10{,}000 \times 64 = 640{,}000 \text{ float values} \\
    \text{Total Bytes} &= 640{,}000 \times 4 \text{ bytes} = 2{,}560{,}000 \text{ bytes} \\
    \text{Kilobytes} &= \frac{2{,}560{,}000}{1024} = \mathbf{2{,}500 \text{ KB}} \approx \mathbf{2.44 \text{ MB}}
    \end{align*}
    At $2.44\text{ MB}$, this embedding matrix fits comfortably in external Flash storage, but exceeds the typical internal SRAM of an ESP32 ($512\text{ KB}$), requiring streaming lookup from Flash or PSRAM.

    \item \textbf{Catastrophic Forgetting Risk}:
    Pretrained Transformer models (like BERT) have already adjusted millions of internal attention weights across billions of sentences into a finely tuned, highly delicate manifold representing syntax, grammar, and world semantics. 
    A large learning rate like $\eta = 0.01$ causes colossal gradient update steps that violently overwrite these pre-existing weights, completely erasing the model's linguistic foundation in the first few batches. A tiny learning rate like $\eta = 0.00002$ gently steers the existing representations toward the specific domain task without breaking general language understanding.

    \item \textbf{Stop Word Pitfall in IoT}:
    Consider the alert: \code{"Valve not closed"}.
    Standard stop-word removal algorithms strip the word \code{"not"}, reducing the text to:
    \code{"Valve closed"}.
    The processed text asserts the exact physical opposite of the true state, blinding the automated safety system to an open valve and potentially causing an industrial tank overflow!
\end{enumerate}
